\documentclass[12pt,a4paper]{article}
\usepackage[utf8]{inputenc}
\usepackage{geometry}
\usepackage{hyperref}
\usepackage{xcolor}
\usepackage{enumitem}
\usepackage{titlesec}

\geometry{
    left=25mm,
    right=25mm,
    top=25mm,
    bottom=25mm
}

\titleformat{\section}
{\normalfont\Large\bfseries\color{purple}}{\thesection}{1em}{}

\titleformat{\subsection}
{\normalfont\large\bfseries\color{blue}}{\thesubsection}{1em}{}

\title{\textbf{Diabetes Risk Assessment App}\\Your Health Companion}
\author{}
\date{}

\begin{document}

\maketitle

\section{Welcome to Your Personal Health Journey}

This mobile app is here to help you understand and manage your diabetes risk through simple, everyday health checks. Think of it as your pocket health companion that guides you through understanding your body's signals and connects you with the right care when you need it.

\section{Getting Started with Your Account}

When you first open the app, you'll see a friendly login screen asking for your email and password. If you're new here, no worries! Just tap the registration button to create your account. You'll only need to provide your name, email address, and choose a secure password. All your information stays private and secure on your device.

\section{Your Personal Dashboard}

Once you're logged in, you'll land on your personal dashboard -- your health command center. Here you'll find everything organized in easy-to-tap cards:

\begin{itemize}
    \item \textbf{Risk Assessment} -- Check your diabetes risk
    \item \textbf{Food Guide} -- Learn about healthy eating
    \item \textbf{Exercise Tips} -- Get moving with physical activity ideas
    \item \textbf{Hospitals} -- Find diabetes care centers near you
    \item \textbf{Doctors} -- Connect with specialist physicians
    \item \textbf{Task List} -- Keep track of your health goals
    \item \textbf{Profile} -- View and manage your account
    \item \textbf{Settings} -- Customize your app experience
\end{itemize}

The dashboard also includes a handy task manager where you can create, edit, and track your personal health-related tasks -- like remembering to take your morning walk or schedule a doctor's appointment.

\section{Understanding Your Diabetes Risk}

The heart of this app is the risk assessment tool. It's like having a conversation about your health -- you'll answer questions about how you've been feeling, and the app helps you understand what those feelings might mean.

You'll go through a simple form that asks about:
\begin{itemize}
    \item Your age and gender
    \item Common health signs like excessive thirst, frequent urination, unexplained weight loss
    \item How you're feeling physically -- any weakness, vision changes, or other symptoms
\end{itemize}

The questions use easy dropdown menus and simple toggle switches, making it quick to complete even on your phone.

\section{Your Personal Risk Results}

After you complete the assessment, the app processes your answers and gives you clear, easy-to-understand results. Your risk level falls into one of five categories:

\begin{itemize}
    \item \textbf{\textcolor{green!60!black}{Very Low Risk}} (Green) -- You're doing great!
    \item \textbf{\textcolor{green!60!black}{Low Risk}} (Green) -- Keep up the good habits
    \item \textbf{\textcolor{yellow!80!black}{Moderate Risk}} (Yellow) -- Time to pay a bit more attention
    \item \textbf{\textcolor{red}{High Risk}} (Red) -- Consider consulting a healthcare provider
    \item \textbf{\textcolor{red}{Critical Risk}} (Red) -- Please seek medical attention
\end{itemize}

The results are shown on color-coded cards so you can instantly understand where you stand. Along with your risk level, you'll get personalized recommendations tailored specifically to your situation.

\section{How Well Does It Work?}

We tested this app with real people to see how helpful it really is. Out of 42 people who tried it, the app's predictions matched their actual medical diagnoses in 33 cases -- that's about 78\% accuracy.

Now, 78\% might not sound perfect, but here's the thing: diabetes can be tricky. Many people don't show any symptoms in the early stages, and symptoms can vary a lot from person to person. This app is designed as a first step -- a screening tool that helps you understand when it might be time to see a doctor, not a replacement for professional medical advice.

\section{Eating Well with Bangladeshi Flavors}

The food guide section is all about eating right while enjoying the foods you love. All the recommendations are written in Bangla and tailored to Bangladeshi cooking and culture. You'll find:

\begin{itemize}
    \item General tips for keeping your blood sugar stable
    \item Lists of foods that are great for you
    \item Foods you might want to enjoy in moderation
    \item Practical advice for everyday meal planning
\end{itemize}

It's not about giving up your favorite foods -- it's about making smart choices that keep you healthy and satisfied.

\section{Getting Active Your Way}

The exercise section makes physical activity approachable and doable. Written in Bangla, it covers:

\begin{itemize}
    \item Aerobic exercises that get your heart pumping
    \item Strength training to build muscle
    \item Clear instructions on proper form and technique
    \item How to exercise safely at your own pace
\end{itemize}

Whether you're a complete beginner or looking to step up your routine, you'll find guidance that works for your fitness level.

\section{Finding Healthcare When You Need It}

The hospital directory is like having a local healthcare map in your pocket. It lists diabetes-specialized hospitals and facilities across Bangladesh, complete with:

\begin{itemize}
    \item What services each facility offers
    \item Where they're located
    \item How to contact them
    \item Emergency service information
\end{itemize}

You can quickly find nearby facilities and get directions, making it easier to get the care you need without the hassle.

\section{Connecting with the Right Doctors}

The physician section helps you find the right specialist for your needs. Each doctor's profile includes:

\begin{itemize}
    \item Their qualifications and training
    \item What they specialize in
    \item Which hospitals they're affiliated with
    \item How to reach them
\end{itemize}

This is especially helpful if you're in a remote area -- you can research doctors before traveling, saving time and ensuring you connect with the right healthcare provider.

\section{Designed for Your Phone, Designed for You}

Everything about this app was built with mobile use in mind:

\begin{itemize}
    \item Big, easy-to-tap buttons perfect for touchscreens
    \item Clean, organized layouts that work on small screens
    \item Fast performance even on older phones
    \item The ability to use key features even without internet
    \item All content in Bangla to eliminate language barriers
\end{itemize}

We know that technology should work for you, not the other way around. That's why this app prioritizes simplicity, accessibility, and genuine usefulness in real-world situations.

\vspace{1cm}
\hrule
\vspace{0.5cm}

\textit{\small Remember: This app is a tool to help you understand your health better, but it's not a substitute for professional medical advice. Always consult with healthcare providers for diagnosis and treatment.}

\end{document}